\documentclass[10pt,a4paper]{amsart}
\usepackage[margin=19mm]{geometry}
\usepackage{amsmath,amssymb,mathtools}
\usepackage[scr=rsfs,cal=euler]{mathalfa}
\usepackage{xfrac}
\usepackage[unicode,hidelinks,bookmarks=false]{hyperref}
\newcommand{\ie}{i.e.\ }
\newcommand{\cf}{cf.\ }
\newcommand{\resp}{respectively\ }
\newcommand{\Subsection}{\subsection}
\providecommand{\nobreakdash}{\nobreak\hbox{-}}
\setlength{\emergencystretch}{2em}
\multlinegap0pt
\usepackage{siunitx}
\usepackage{siunitx}
\usepackage{siunitx}
\usepackage{siunitx}
\usepackage{amssymb}
\usepackage{bm}
\usepackage{cancel}
\usepackage{enumerate}
\usepackage[shortlabels]{enumitem}
\usepackage{mathtools}
\usepackage{siunitx}
\usepackage{tikz}
\newtheorem{theorem}{Theorem}
\newtheorem{lemma}{Lemma}
\newtheorem{corollary}{Corollary}
\newtheorem{proposition}{Proposition}
\usepackage{cleveref}
\crefname{theorem}{Theorem}{Theorems}
\crefname{lemma}{Lemma}{Lemmas}
\crefname{corollary}{Corollary}{Corollarys}
\crefname{proposition}{Proposition}{Propositions}
\newcommand{\Del}{\Delta}
\DeclareMathOperator{\Dil}{Dil}
\newcommand{\G}{\Gamma}
\newcommand{\Z}{\mathbb{Z}}
\newcommand{\al}{\alpha}
\newcommand{\calD}{\mathcal{D}}
\newcommand{\del}{\delta}
\renewcommand{\k}{\kappa}
\newcommand{\lam}{\lambda}
\newcommand{\m}{\mu}
\newcommand{\scrD}{\mathscr{D}}
\title{A geometric correspondence for reparameterizations of geodesic flows}
\author{Stephen Cantrell, D\'idac Mart\'inez-Granado, Eduardo Reyes}
\date{Source: arXiv:2605.02585v2 (CC BY 4.0)}
\begin{document}
\maketitle

\noindent\emph{Source and licence.} A geometric correspondence for reparameterizations of geodesic flows. Version arXiv:2605.02585v2 (CC BY 4.0), distributed by arXiv under the Creative Commons Attribution 4.0 International licence, http://creativecommons.org/licenses/by/4.0/, as recorded in the arXiv licence metadata for this exact version and shown on the abstract page https://arxiv.org/abs/2605.02585v2. Full licence notice: \texttt{licenses/arxiv-2605.02585-CC-BY-4.0.txt}.\par\smallskip

\noindent\textit{Editorial scope.} Counted unit: the article's theorem thm.greendense (source0.tex lines 2512--2514). The article's complete original proof of it is delivered with it (source0.tex lines 2638--2661, one \texttt{proof} environment of 24 lines, which the article places away from the statement). No mathematics is written, repaired or re-derived: the statement and the proof are the article's own lines, in the article's own notation, and no retained line is substituted or re-flowed.  Retained uncounted as the source-local context the proof needs: lines 2544--2549; lines 2570--2573; lines 2593--2598; lines 2611--2615.  Nothing was omitted from any retained span, so the package carries no omission record.  No retained line contains a citation, so the wrapper carries no bibliography and no bibliography entry.  No retained line contains a figure environment, an image-inclusion command or drawing code, so no image is delivered and the package declares no asset dependency.  Editorial formatting only, in addition: an AMS \texttt{amsart} preamble that restates the article's own declarations and macros used by the retained lines, namely the environments \texttt{theorem}, \texttt{lemma}, \texttt{corollary}, \texttt{proposition} on separate counters, together with the standard packages those lines need.  The retained environments print in this document as Theorem 1, Lemma 1, Corollary 1, Proposition 1 in order of appearance, whereas the article prints the same items with the numbering of its own full document.  The complete native source of the article as distributed in the arXiv e-print for this exact version is bundled unchanged as \texttt{sources/199727/source0.tex}.
\begin{lemma}\label{lem.worddensesphere}
    Let $\G$ be any group and let $d$ be an unbounded left-invariant $\al$-roughly geodesic (non-necessarily symmetric) metric on $\G$. Then for any $l>\del>2(\lceil\al \rceil+\al+2)$ the set $S_l=\{g\in \G: l-\del<d(o,g)\leq l\}$ generates $\G$ as a semi-group and there exists $M=M_{l,\al}>0$ such that
\begin{equation}\label{eq.WMdensesphere}
        (l-\lceil{\al \rceil})|g|_{S_l}-M\leq d(o,g)\leq l|x|_{S_l} \quad \text{ for all }g\in \G.
    \end{equation}
\end{lemma}

\begin{corollary}\label{coro.convEGR}
    Let $\G$ be non-elementary hyperbolic and $d\in \calD_\G$ be $\al$-roughly geodesic. Then there exists a constant $\beta>0$ depending only on $d$ such that for any $l>\del >2(\lceil \al \rceil+\al+2)$ we have
    \[l^{-1}(\log(\# S_l)-\beta)\leq v_d\leq (l-\lceil \al\rceil)^{-1}\log(\# S_l).\]
\end{corollary}

\begin{lemma}\label{lem.greeneasyineq}
Let $\G$ be an infinite finitely generated group that is not virtually $\Z$ or $\Z^2$ and let $S\subset \G$ be a finite set generating $\G$ as a semi-group. If $\m=\m_S$ denotes the uniform probability measure supported on $S$, then for all $g\in \G$ we have
\begin{equation*}
    d_\mu(o,g)\leq \log\left(\#S\right)|g|_S+\log\left(G_\m(o,o)(1-(\#S)^{-1})\right).
\end{equation*}
\end{lemma}

\begin{proposition}\label{prop.greenmetricineqHYP}
    Let $\G$ be a non-elementary hyperbolic group and consider $d\in \calD_\G$. Then there exists $\del_0$ satisfying the following. Let $\del>\del_0$ and for $l$ large enough define $S_l=S_l^\del:=\{g\in \G \colon l-\del< d(o,g)\leq l\}$. Then $S_l$ generates $\G$ (as a group) and if $\m_l$ is the uniform probability measure supported on $S_l$, then there exists $\k_l>0$ such that for all $g\in \G$ we have \begin{equation}\label{eq.propdGreen}
            d_{\m_l}(o,g)\geq \frac{v_dd(o,g)}{2}+|g|_{S_l}\cdot \left(\log\left(\frac{\#S_l}{D_0\left(\frac{l}{\del}+1\right)}\right)-\frac{v_d l}{2}\right)-\k_l.
        \end{equation}
\end{proposition}

\begin{theorem}[Density of Green metrics]\label{thm.greendense}
    Let $\G$ be a non-elementary hyperbolic group and $d\in \calD_\G$. Then there exists a constant $\del_1>0$ satisfying the following. Let $\del>\del_1$ and $\mu_l$ be the uniform probability measure supported on $S_l=\{g\in \G \colon l-\del<d(o,g)\leq l\}$. Then $S_l$ generates $\G$ for $l$ large enough and $[d_{\mu_l}]$ converges to $[d]$ in $\scrD_\G$.
\end{theorem}

\begin{proof}[Proof of \Cref{thm.greendense}]
    Let $d\in \calD_\G$, which we assume to be $\al$-roughly geodesic. We fix $\del>\max(\del_0,2(\lceil \al\rceil+\al+2)),$ so that both \Cref{lem.worddensesphere} and  \Cref{prop.greenmetricineqHYP} apply. For $l>0$ large enough and $g\in \G$ this gives us 
    \begin{align*}
                d_{\m_l}(o,g) &\geq \frac{v_dd(o,g)}{2}+|g|_{S_l}\cdot \left(\log\left(\frac{\#S_l}{C\left(\frac{l}{\del}+1\right)}\right)-\frac{v_d l}{2}\right)-\k_l\\
                & \geq  \frac{v_dd(o,g)}{2}+d(o,g)\cdot \left(\frac{\log(\# S_l)}{l}-\frac{\log(C\left(\frac{l}{\del}+1\right))}{l}-\frac{v_d}{2}\right)-\k_l\\
                & = d(o,g)\cdot \left(\frac{\log(\# S_l)}{l}-\frac{\log(C\left(\frac{l}{\del}+1\right))}{l}\right)-\k_l.
    \end{align*}    
In particular, we obtain
    \begin{equation}\label{eq.ineqDil1}
        \Dil(d,d_{\m_l})\leq \left(\frac{\log{\# S_l}}{l}-\frac{\log(C(\frac{l}{\del}+1))}{l}\right)^{-1},
    \end{equation}
 and the right hand side of \Cref{eq.ineqDil1} tends to $v_d^{-1}$ as $l$ tends to  $\infty$ by \Cref{coro.convEGR}.

 Also, applying \Cref{lem.greeneasyineq} to $S=S_l$ and \Cref{lem.worddensesphere} we get
\begin{align*}
    d_{\m_l}(o,g)& \leq \log(\# S_l)|g|_{S_l}+\log(G_{\m_l}(o,o)^{-1}(1-(\#S_l)^{-1})) \\
    & \leq \frac{\log(\# S_l)}{(l-\lceil \al\rceil)}\cdot d(o,g)+\log(G_{\m_l}(o,o)^{-1}(1-(\#S_l)^{-1}))+\frac{M}{(l-\lceil \al \rceil)}, 
\end{align*}
and hence
\begin{equation}\label{eq.greenDil2}
    \Dil(d_{\m_l},d)\leq  \frac{\log(\# S_l)}{(l-\lceil \al\rceil)}.
\end{equation}
 The right hand side of \Cref{eq.greenDil2} tends to $v_d$ as $l$ tends to infinity by \Cref{coro.convEGR}, implying that  $\Del([d_{\lam_l}],[d])=\log(\Dil(d_{\lam_l},d)\Dil(d,d_{\lam_l}))$ tends to $0$ as $l$ tends to infinity, as desired.
\end{proof}

\end{document}
